% !TeX root = main.tex
% Synchronized with 00_摘要.md v0.5; Tsai format.
\begin{abstractzh}
\setcounter{page}{1}

移動機械臂結合底盤移動與手臂操作能力，可執行室內搬運與關節式物件操作。然而，動態障礙、操作可達性與手臂追蹤落後，增加導航與協同控制的難度。本研究探討動態環境下移動機械臂的預測導航與受約束操作，建立結合導航、全身協調與視覺導引的控制流程。

導航採用以 \(SE(2)\) 誤差描述的幾何模型預測控制，結合預測式控制屏障條件及原始雷射掃描保護層；協同控制將手臂位置設定點與追蹤動態納入預測，共同最佳化底盤與手臂命令。

模擬結果顯示，所提出方法在三種方法共同可評估的 83 條導航路線中全部到達，其中 1 條發生接觸。三對抽屜完整策略比較顯示，移動中協同操作相對嚴格停車策略縮短任務時間 13.38--14.05\%，兩組均通過原任務判準。另三對實驗顯示，五步預測的整控制窗末端位置誤差與手臂命令反轉率低於單步預測。視覺單瓶抓放與理想抓持門操作已取得指定案例成果，門操作仍採非即時計算；完整住宅取物、全身動態避障及實機即時執行尚待驗證。

\keywordzh{\mbox{移動機械臂}、\allowbreak\mbox{幾何模型預測控制}、\allowbreak\mbox{控制屏障函數}、\allowbreak\mbox{全身協調}、\allowbreak\mbox{移動中操作}、\allowbreak\mbox{視覺導引}}
\end{abstractzh}
